\subsection{One coefficient for both limiting components}
\label{sec:joint-coefficient}
We spell out the coefficient selection used in the preceding realization.
All decompositions in this subsection are under the fixed auxiliary
probability $Q$, on the original price $S$.
Write $S-S_0=M^S+A^S$, with zero-start canonical components.
Let $H_n$ be the actual approximating coefficients and write
\[
 M_n=H_n\cdot M^S,\qquad A_n=H_n\cdot A^S.
\]
The preceding component-series construction supplies regular limits $M,A$,
locally uniform convergence of $M_n$ to $M$, variation convergence of $A_n$
to $A$, and, almost surely on each finite horizon $T$,
\[
 \sum_n\Var_{[0,T]}(A_{n+1}-A_n)<\infty.
\]
It also supplies common joint stops on each of which the martingale
terminals are bounded in $L^2(Q)$, uniformly in the approximation index.

\begin{lemma}[Simultaneous coefficient selection]\label{lem:joint-coefficient}
Under these conditions there is one predictable real coefficient $H$ whose
completed martingale and finite-variation integrals on every joint coordinate
are the corresponding stops of $M$ and $A$. No absolute continuity between
the energy measure and the variation measure is required.
\end{lemma}
\begin{proof}
Fix first a deterministic $T>0$. Write $v_T(\omega)=|d(A^S)^T|(\omega)$ for
the source's pathwise variation measure. Choose a finite measurable random
variable $R_T\geq0$ dominating
\[
 v_T(\omega)([0,\infty))+\Var_{[0,T]}A_0+
 \sum_n\Var_{[0,T]}(A_{n+1}-A_n),
\]
Such a variable exists by the pathwise summability already proved.
The construction may add the summable martingale envelopes to $R_T$;
this makes the same weight usable for all component estimates.
Set $a_T=(1+R_T)^{-2}$, replacing it by $1$ on the null exceptional set.
It is strictly positive almost surely. Define a finite measure on the
predictable sigma-field by
\[
 \nu_T(B)=E_Q\left[a_T\int1_B(t,\omega)\,dv_T(\omega)(t)\right].
\]
This changes only the variation control; martingale and energy estimates
remain under $Q$.

For every approximating integral the pathwise density identity gives
\[
 \begin{aligned}
 \int|H_0|\,d\nu_T&=E_Q[a_T\Var_{[0,T]}A_0]<\infty,\\
 \sum_n\int|H_{n+1}-H_n|\,d\nu_T
 &=E_Q\left[a_T\sum_n\Var_{[0,T]}(A_{n+1}-A_n)\right]<\infty.
 \end{aligned}
\]
Indeed $d(A_{n+1}-A_n)=(H_{n+1}-H_n)dA^S$ and the variation of a signed
measure with density $f$ is $|f|$ times its variation measure.
Tonelli and completeness of $L^1$ imply that the full coefficient sequence
has a limit $h_T$, both $\nu_T$-almost everywhere and in $L^1(\nu_T)$.
Equivalently, outside a $Q$-null set it converges $v_T(\omega)$-almost
everywhere and in pathwise $L^1(v_T(\omega))$: its successive pathwise
$L^1$ distances are summable and its initial term is integrable.

For the countably many joint terminal coordinates choose $b_r>0$ with
$\sup_n\|M_n(\rho_r)\|_2\leq b_r$. The vectors
$(2^{-(r+1)/2}M_n(\rho_r)/b_r)_r$ are bounded in the Hilbert direct sum
of the terminal $L^2(Q)$ spaces. Weak subsequential compactness and Mazur's
lemma give forward convex averages converging strongly in this direct sum.
Transfer the subsequence weights back to the original indices; their support
still tends forward. Thus one family $w_{n,k}$ makes every terminal
coordinate Cauchy. Define $\bar H_n=\sum_kw_{n,k}H_k$, and use the same
weights for $\bar M_n$.
Let $q_r$ be the predictable energy measure on the $r$th coordinate.
The completed isometry gives
\[
 \|\bar H_n-\bar H_m\|_{L^2(q_r)}
 =\|\bar M_n(\rho_r)-\bar M_m(\rho_r)\|_{L^2(Q)}.
\]
Here $\rho_r$ is the finite coordinate stop, including its deterministic
horizon, and $\bar M_n$ uses precisely the same weights.
Thus these coefficients are Cauchy in each of the countably many
$L^2(q_r)$ spaces. Choose one strictly increasing $f$ with
\[
 \|\bar H_{f(j+1)}-\bar H_{f(j)}\|_{L^2(q_r)}\leq2^{-j}
 \quad(r\leq j).
\]
Each $q_r$ is finite, so Cauchy--Schwarz and Tonelli give pointwise convergence
$q_r$-almost everywhere. Define $H$ to be the pointwise limit where this
sequence converges in $\mathbb R$, and zero elsewhere. Its convergence set
is predictable, so $H$ is predictable; $L^2$ completeness gives
$\bar H_{f(j)}\to H$ in $L^2(q_r)$ for every $r$.

The same definition determines the variation component, including regions
where $q_r$ vanishes. At each point where $H_n\to h_T$, forward convexity gives
\[
 |\bar H_n-h_T|\leq\sup_{k\geq n}|H_k-h_T|\longrightarrow0.
\]
Hence $H=h_T$ $\nu_T$-almost everywhere, and also
$v_T(\omega)$-almost everywhere for almost every $\omega$.
In particular $H\in L^1(\nu_T)$ and is pathwise integrable against $v_T$.
A joint source prefix is a further stop before its deterministic horizon,
so its variation measure is a restriction of $v_T$. The same integrability
therefore holds on that coordinate. This explicitly supplies both
$L^2(q_r)$ and $L^1$ membership for one coefficient.

The completed martingale isometry and Doob's inequality show that the
martingale integrals of $\bar H_{f(j)}$ converge to that of $H$ on coordinate
$r$. Forward averages preserve the already existing limit $M$, and uniqueness
of limits in probability identifies this integral with $M^{\rho_r}$.
On the variation side, for almost every path,
\[
 \Var_{[0,T]}\bigl(H_n\cdot A^S-H\cdot A^S\bigr)
 =\int|H_n-H|\,dv_T\longrightarrow0.
\]
The left-hand integrals also converge in variation to $A^T$.
Zero initial values and uniqueness of limits identify $H\cdot(A^S)^T=A^T$.
Restricting these signed-measure identities to $(0,\rho_r]$ proves the
finite-variation identification on the joint coordinate, with the endpoint
jump included. Countably many coordinates and right continuity give one
null set for all process equalities. Thus the same $H$ realizes both parts.
\end{proof}
The coefficients in different outer stopped stages need not agree.
Section~\ref{sec:integral-calculus} pastes them on disjoint stopping intervals after
return to the original measure. Lemma~\ref{lem:joint-coefficient} supplies
the stopped realizations required there; it does not assume integral-range
closedness in selecting $H$.
